\preface

This book is written as a modern, research-grounded introduction to large language models. The prose, organization, notation, examples, and arguments are original. Each chapter is intended to connect three views that are too often taught separately: the mathematical object being optimized, the system that makes the optimization feasible, and the evaluation procedure that determines whether the result is useful.

The book deliberately avoids an implementation-by-implementation narrative. Instead, it follows the lifecycle of a model: data and tokenization, architecture, pretraining, distributed systems, serving, adaptation, alignment, retrieval, reasoning, multimodality, and safety. Historical systems such as GPT, Alpaca, LoRA, RLHF, and LLaMA are treated as conceptual anchors rather than as recipes to copy without context.

The exposition is evidence-first. Claims about algorithms, scaling behavior, evaluation, or deployment are tied to primary papers or technical reports whenever possible. No third-party slide text, media transcript, dataset record, or code listing is reproduced as manuscript prose.
